\section{Atomic Value Movement}
\label{sec:cross-chain}

Trading on the D-Chain requires moving value into and out of it from the public Lux
network. This is the responsibility of the atomic proxy (Section~\ref{sec:architecture}),
and it reuses the mechanism the Lux primary network already uses to exchange assets
between the X- and C-Chains: atomic shared-memory import/export, committed at block
acceptance. There is no external relayer trusted to move funds and no bridge holding
custody outside consensus.

\subsection{Shared-memory import/export}

A transfer between a public-network chain and the D-Chain is expressed as a matched
export/import pair applied atomically at \texttt{Block.Accept}:

\begin{algorithm}[H]
\caption{Atomic import into the D-Chain}
\begin{algorithmic}[1]
\STATE \textbf{Export} on the source chain: lock value and record an export to the
D-Chain's shared-memory address space.
\STATE \textbf{Relay}: the proxy carries the order and the pending import to the D-Chain.
\STATE \textbf{Apply at Accept}: when the D-Chain block is accepted, the matched
import/export is applied to shared memory in one commit---both legs take effect or
neither does.
\STATE \textbf{Settle}: fills from matching settle against the imported balance on the
D-Chain; the symmetric export path returns value to the source chain the same way.
\end{algorithmic}
\end{algorithm}

Because the import and export are applied in a single shared-memory commit, value
conservation is structural: the proxy can lose its liveness (fail to relay) but it cannot
create or destroy value, and a block that is verified and then rejected applies nothing,
stranding no funds. The conservation property is exercised by an end-to-end test that
drives a transfer across the public chain, the proxy, and the D-Chain and asserts that
total balances are preserved all-or-nothing.

\subsection{Native asset convention}

For EVM compatibility the facade exposes native LUX as \texttt{address(0)} rather than a
wrapped token, matching the Uniswap-V4 native-currency convention. Wrapped liquidity, if
present, is a distinct asset---never the canonical key for native LUX. Internally the
canonical asset resolves to the D-Chain's native identity, so there is one representation
of native value on the matching path.

\subsection{Settlement cadence versus matching cadence}

Atomic value movement runs at block cadence: a transfer settles when a D-Chain block is
accepted, not per order. This is deliberate. The latency-critical leg is matching, which
is co-located inside the venue and measured in nanoseconds per fill
(Section~\ref{sec:gpu-matching}); the cross-chain settlement leg is the proxy's
responsibility and runs once per block. Separating the two---fast matching on-host, atomic
settlement at block cadence---keeps the hot path off the network relay, whose round-trip
would otherwise dominate per-order latency.
